\documentclass[11pt]{article}

\usepackage[margin=1in]{geometry}
\usepackage{amsmath,amssymb}
\usepackage[T1]{fontenc}
\usepackage{lmodern}
\usepackage{hyperref}
\usepackage{array}
\usepackage{booktabs}

\title{Full Mathematical Note on the Fixed-\texorpdfstring{$y$}{y} Quantum SCF Critical-Point Solver}
\author{}
\date{\today}

\begin{document}
\maketitle

\begin{abstract}
This document gives a fully mathematical explanation of the fixed-$y$ quantum
critical-point solver used in \texttt{nuclear\_matter\_workflows}. It explains
the finite-temperature physics, the meaning of the fitted Clausius parameters,
the inner self-consistent inversion for the effective chemical potentials
$\mu_p^\ast$ and $\mu_n^\ast$, and the outer critical-point solve for
$T_c$ and $n_c$. A worked numerical example using the repository output for
$K_0=280\ \mathrm{MeV}$ is included.
\end{abstract}

\tableofcontents

\section{What This Solver Takes as Input}

The finite-temperature fixed-$y$ solver does \emph{not} fit the interaction
parameters. It starts from a row already produced by the zero-temperature
asymmetric Clausius fit:
\[
a_n,\qquad a_{pn},\qquad b_n,\qquad b_{pn},\qquad c.
\]

So the logical order is
\[
\text{zero-temperature asymmetric fit}
\quad \Longrightarrow \quad
\text{finite-temperature fixed-}y\text{ critical solve}.
\]

The role of the fixed-$y$ solver is therefore:
\begin{quote}
Given one calibrated asymmetric Clausius interaction and one chosen proton
fraction $y$, find the finite-temperature liquid-gas critical point.
\end{quote}

\section{What Is the Critical Point?}

At fixed composition $y$, the solver builds the pressure as a function of
temperature and baryon density,
\[
P(T,n;y).
\]

The liquid-gas critical point is defined by the standard inflection-point
conditions
\begin{align}
\left(\frac{\partial P}{\partial n}\right)_T &= 0, \\
\left(\frac{\partial^2 P}{\partial n^2}\right)_T &= 0.
\end{align}

This is the same mathematics as the textbook van der Waals critical point.

Physically:
\begin{itemize}
\item below $T_c$, the isotherm can have a liquid-gas coexistence structure;
\item at $T_c$, the coexistence endpoint is reached;
\item above $T_c$, there is no first-order liquid-gas transition.
\end{itemize}

\section{Step 1: Species Densities at Fixed Proton Fraction}

The solver works at fixed proton fraction $y$, so
\begin{align}
n_p &= yn, \\
n_n &= (1-y)n.
\end{align}

Only the range
\[
0 \le y \le 0.5
\]
needs to be considered, because interchanging protons and neutrons maps
$y \leftrightarrow 1-y$.

\section{Step 2: Species-Dependent Excluded Volume}

The asymmetric Clausius model introduces different excluded-volume couplings for
same-species and proton-neutron channels:
\begin{align}
x_p &= b_n n_p + b_{pn} n_n, \\
x_n &= b_{pn} n_p + b_n n_n.
\end{align}

The free-volume fractions are then
\begin{align}
f_p &= 1-x_p, \\
f_n &= 1-x_n.
\end{align}

These must satisfy
\[
f_p>0,\qquad f_n>0.
\]
Otherwise, the state is unphysical and the code rejects it.

\subsection{Physical Meaning}

This is the same excluded-volume idea as at $T=0$, but now applied separately
to protons and neutrons.

An easy analogy is a two-component gas of red and blue hard balls:
\begin{itemize}
\item a red ball excludes some volume from other red balls;
\item a red ball may exclude a different effective volume from blue balls.
\end{itemize}

That is why the model distinguishes $b_n$ and $b_{pn}$.

\section{Step 3: Target Ideal Densities}

Because only fractions $f_p$ and $f_n$ of the volume are available, the
corresponding ideal-gas target densities are
\begin{align}
n_{p,\mathrm{id}}^{\mathrm{target}} &= \frac{n_p}{f_p}, \\
n_{n,\mathrm{id}}^{\mathrm{target}} &= \frac{n_n}{f_n}.
\end{align}

These are the densities that the finite-temperature proton and neutron Fermi
gases must reproduce.

\section{Step 4: Finite-Temperature Quantum Fermi Gas}

\subsection{Fermi-Dirac Distribution}

At nonzero temperature, the occupation factor is
\begin{equation}
f_{\mathrm{FD}}(E;\mu,T)=\frac{1}{e^{(E-\mu)/T}+1}.
\end{equation}

This replaces the sharp $T=0$ Fermi sphere.

\subsection{Relativistic Dispersion Relation}

The one-particle energy is
\begin{equation}
E(k)=\sqrt{(\hbar c\,k)^2+m_N^2}.
\end{equation}

\subsection{Ideal Density at Finite Temperature}

For degeneracy $d$,
\begin{equation}
n_{\mathrm{id}}(T,\mu)
=
\frac{d}{2\pi^2}\int_0^\infty k^2 f_{\mathrm{FD}}(E;\mu,T)\,dk.
\end{equation}

\subsection{Ideal Pressure at Finite Temperature}

The ideal pressure is
\begin{equation}
p_{\mathrm{id}}(T,\mu)
=
\frac{d}{6\pi^2}\int_0^\infty
\frac{(\hbar c)^2k^4}{E(k)}f_{\mathrm{FD}}(E;\mu,T)\,dk.
\end{equation}

\subsection{Why the Solver Is Called Quantum}

The finite-temperature stage is called quantum because it keeps the full
Fermi-Dirac statistics. It is not using a classical Boltzmann gas.

\subsection{Species Degeneracy}

In the fixed-$y$ solver, protons and neutrons are treated separately, so each
species uses
\[
d=2
\]
rather than $d=4$.

\section{Step 5: The Inner SCF Problem}

At a trial thermodynamic point $(T,n)$, the solver must determine effective
chemical potentials $\mu_p^\ast$ and $\mu_n^\ast$ such that
\begin{align}
n_{\mathrm{id}}(T,\mu_p^\ast)
&=
n_{p,\mathrm{id}}^{\mathrm{target}}, \\
n_{\mathrm{id}}(T,\mu_n^\ast)
&=
n_{n,\mathrm{id}}^{\mathrm{target}}.
\end{align}

This is the first self-consistent problem.

\subsection{Write It as Root Finding}

For one species, define
\begin{equation}
F(\mu^\ast)=n_{\mathrm{id}}(T,\mu^\ast)-n_{\mathrm{id}}^{\mathrm{target}}.
\end{equation}

The problem is to solve
\begin{equation}
F(\mu^\ast)=0.
\end{equation}

\subsection{Initial Guess}

The code starts from a zero-temperature estimate:
\begin{align}
k_F &= \left(\frac{6\pi^2 n_{\mathrm{id}}^{\mathrm{target}}}{d}\right)^{1/3}, \\
\mu_{\mathrm{seed}} &= \sqrt{(\hbar c\,k_F)^2+m_N^2}.
\end{align}

This is a natural starting point because at low temperature the finite-$T$
chemical potential is close to the $T=0$ Fermi energy.

\subsection{Newton-Like Update}

The code estimates the local derivative numerically:
\begin{equation}
\chi(\mu)
\approx
\frac{n_{\mathrm{id}}(T,\mu+\delta\mu)-n_{\mathrm{id}}(T,\mu-\delta\mu)}
{2\delta\mu}.
\end{equation}

Then the Newton-like update is
\begin{equation}
\mu_{\mathrm{new}}
=
\mu-\lambda_\mu \frac{F(\mu)}{\chi(\mu)},
\end{equation}
where $\lambda_\mu$ is a damping factor.

\subsection{Why Damping Is Used}

Pure Newton iteration can overshoot if the function is stiff or the initial
guess is not close enough. Damping makes the method more stable.

An easy analogy is walking downhill with shorter steps so that you do not jump
past the valley.

\section{Step 6: The Fixed-\texorpdfstring{$y$}{y} Pressure}

Once $\mu_p^\ast$ and $\mu_n^\ast$ are known, the proton and neutron ideal
pressures are
\begin{align}
p_p^{\mathrm{id}} &= p_{\mathrm{id}}(T,\mu_p^\ast), \\
p_n^{\mathrm{id}} &= p_{\mathrm{id}}(T,\mu_n^\ast).
\end{align}

\subsection{Effective Attractive Coefficient}

The asymmetric attractive coefficient at fixed $y$ is
\begin{equation}
C(y)
=
a_n\big[y^2+(1-y)^2\big]+2a_{pn}y(1-y).
\end{equation}

This is simply the pair-counting structure for the two-component system:
\begin{itemize}
\item $y^2$ weights proton-proton pairs;
\item $(1-y)^2$ weights neutron-neutron pairs;
\item $2y(1-y)$ weights proton-neutron pairs.
\end{itemize}

\subsection{Interaction Pressure}

The mean-field interaction pressure is
\begin{equation}
P_{\mathrm{int}}(n,y)
=
-\frac{C(y)n^2}{(1+cn)^2}.
\end{equation}

\subsection{Total Pressure}

Therefore the fixed-$y$ pressure branch is
\begin{equation}
P(T,n;y)
=
p_p^{\mathrm{id}}(T,\mu_p^\ast)
+
p_n^{\mathrm{id}}(T,\mu_n^\ast)
-\frac{C(y)n^2}{(1+cn)^2}.
\end{equation}

\section{Step 7: Why the Symmetric Limit Is Recovered at \texorpdfstring{$y=0.5$}{y=0.5}}

At $y=0.5$,
\[
n_p=n_n=\frac{n}{2}.
\]

Then
\begin{align}
x_p &= b_n\frac{n}{2}+b_{pn}\frac{n}{2}
= \frac{b_n+b_{pn}}{2}n
= b_{\mathrm{avg}}n, \\
x_n &= b_{pn}\frac{n}{2}+b_n\frac{n}{2}
= b_{\mathrm{avg}}n.
\end{align}

So
\[
f_p=f_n=1-b_{\mathrm{avg}}n.
\]

For the attractive coefficient,
\begin{align}
C(0.5)
&=
a_n\left[\left(\frac12\right)^2+\left(\frac12\right)^2\right]
+2a_{pn}\left(\frac12\right)\left(\frac12\right) \\
&=
\frac{a_n+a_{pn}}{2}
=
a_{\mathrm{avg}}.
\end{align}

So the fixed-$y$ solver reduces exactly to the symmetric Clausius branch.

\section{Step 8: The Outer Critical-Point Problem}

The second self-consistent problem is to solve
\begin{align}
G_1(T,n) &= \left(\frac{\partial P}{\partial n}\right)_T = 0, \\
G_2(T,n) &= \left(\frac{\partial^2 P}{\partial n^2}\right)_T = 0.
\end{align}

\subsection{Numerical First Derivative}

The code evaluates
\begin{equation}
\left(\frac{\partial P}{\partial n}\right)_T
\approx
\frac{P(T,n+h_n)-P(T,n-h_n)}{2h_n}.
\end{equation}

\subsection{Numerical Second Derivative}

It also evaluates
\begin{equation}
\left(\frac{\partial^2 P}{\partial n^2}\right)_T
\approx
\frac{P(T,n+h_n)-2P(T,n)+P(T,n-h_n)}{h_n^2}.
\end{equation}

\subsection{Seed Search}

Before iterating, the solver scans a coarse rectangular grid in $(T,n)$ and
computes a normalized score
\begin{equation}
S(T,n)=\left(\frac{G_1}{s_1}\right)^2+\left(\frac{G_2}{s_2}\right)^2,
\end{equation}
where $s_1$ and $s_2$ are characteristic scales obtained from the scan.

The point with the smallest score is used as the initial seed.

\subsection{Outer Newton-Like Update}

The update equations are
\begin{align}
\Delta T
&=
-\lambda_T \frac{G_1}{\partial G_1/\partial T}, \\
\Delta n
&=
-\lambda_n \frac{G_2}{\partial G_2/\partial n},
\end{align}
where the denominator derivatives are also estimated numerically.

The trial step is then damped, clamped to a safe maximum size, and checked
through a backtracking line search.

\subsection{Why Two SCF Layers Exist}

This solver contains two nested iterative layers:
\begin{itemize}
\item the inner one solves $(\mu_p^\ast,\mu_n^\ast)$ at fixed $(T,n)$;
\item the outer one solves $(T_c,n_c)$ using the pressure built from the inner
solution.
\end{itemize}

\section{Worked Example: \texorpdfstring{$K_0=280$}{K0=280}, \texorpdfstring{$y=0.3$}{y=0.3}}

Take the fitted interaction
\begin{align}
a_n &= 250.33652601237475\ \mathrm{MeV\ fm}^3, \\
a_{pn} &= 658.4888689143987\ \mathrm{MeV\ fm}^3, \\
b_n &= 1.417423248321993\ \mathrm{fm}^3, \\
b_{pn} &= 2.455476198790412\ \mathrm{fm}^3, \\
c &= 4.11410812010825\ \mathrm{fm}^3.
\end{align}

For this interaction and $y=0.3$, the solver finds
\begin{align}
T_c &= 14.063530458374572\ \mathrm{MeV}, \\
n_c &= 0.048943415460203074\ \mathrm{fm}^{-3}, \\
P_c &= 0.2110688843825943\ \mathrm{MeV\ fm}^{-3}, \\
\mu_p^\ast &= 934.7468901046402\ \mathrm{MeV}, \\
\mu_n^\ast &= 950.2022651752985\ \mathrm{MeV}.
\end{align}

\subsection{Species Densities}

Compute
\begin{align}
n_p &= yn_c = (0.3)(0.048943415460203074)
= 0.014683024638060921\ \mathrm{fm}^{-3}, \\
n_n &= (1-y)n_c = (0.7)(0.048943415460203074)
= 0.03426039082214215\ \mathrm{fm}^{-3}.
\end{align}

\subsection{Excluded-Volume Factors}

Then
\begin{align}
x_p
&=
b_n n_p + b_{pn} n_n \\
&=
(1.417423248321993)(0.014683024638060921)
+(2.455476198790412)(0.03426039082214215) \\
&=
0.10493763470269968,
\end{align}
and
\begin{align}
x_n
&=
b_{pn} n_p + b_n n_n \\
&=
(2.455476198790412)(0.014683024638060921)
+(1.417423248321993)(0.03426039082214215) \\
&=
0.08461529197291352.
\end{align}

So
\begin{align}
f_p &= 1-x_p = 0.8950623652973003, \\
f_n &= 1-x_n = 0.9153847080270865.
\end{align}

\subsection{Ideal Target Densities}

The inner SCF therefore must solve for chemical potentials such that
\begin{align}
n_{p,\mathrm{id}}^{\mathrm{target}}
&=
\frac{n_p}{f_p}
=
\frac{0.014683024638060921}{0.8950623652973003}
=
0.016404471026087514\ \mathrm{fm}^{-3},
\end{align}
and
\begin{align}
n_{n,\mathrm{id}}^{\mathrm{target}}
&=
\frac{n_n}{f_n}
=
\frac{0.03426039082214215}{0.9153847080270865}
=
0.03742731391699016\ \mathrm{fm}^{-3}.
\end{align}

\subsection{Effective Attractive Coefficient}

At $y=0.3$,
\begin{align}
C(0.3)
&=
a_n(0.3^2+0.7^2)+2a_{pn}(0.3)(0.7) \\
&=
a_n(0.58)+a_{pn}(0.42) \\
&=
(250.33652601237475)(0.58)
+(658.4888689143987)(0.42) \\
&=
421.76051003122484\ \mathrm{MeV\ fm}^3.
\end{align}

\subsection{Interaction Pressure}

The attractive pressure contribution is
\begin{align}
P_{\mathrm{int}}
&=
-\frac{C(0.3)n_c^2}{(1+cn_c)^2} \\
&=
-\frac{(421.76051003122484)(0.048943415460203074)^2}
{\big(1+(4.11410812010825)(0.048943415460203074)\big)^2} \\
&=
-0.7000179980681587\ \mathrm{MeV\ fm}^{-3}.
\end{align}

Since
\[
P_c=0.2110688843825943\ \mathrm{MeV\ fm}^{-3},
\]
the total ideal-gas proton-plus-neutron contribution must be
\begin{align}
p_p^{\mathrm{id}}+p_n^{\mathrm{id}}
&=
P_c-P_{\mathrm{int}} \\
&=
0.2110688843825943-(-0.7000179980681587) \\
&=
0.911086882450753\ \mathrm{MeV\ fm}^{-3}.
\end{align}

So the final critical pressure is the result of a competition between:
\begin{itemize}
\item positive quantum ideal-gas pressure;
\item negative mean-field attraction.
\end{itemize}

\section{Worked Symmetric Check: \texorpdfstring{$K_0=280$}{K0=280}, \texorpdfstring{$y=0.5$}{y=0.5}}

For the same interaction at $y=0.5$, the solver gives
\begin{align}
T_c &= 16.4817165975974\ \mathrm{MeV}, \\
n_c &= 0.0520814359187431\ \mathrm{fm}^{-3}, \\
P_c &= 0.261959538318564\ \mathrm{MeV\ fm}^{-3}.
\end{align}

At $y=0.5$,
\[
n_p=n_n=\frac{n_c}{2}.
\]
Then
\begin{align}
x_p &= x_n = b_{\mathrm{avg}}n_c
=
(1.9364497235562026)(0.0520814359187431) \\
&=
0.10085308218726015,
\end{align}
so
\[
f_p=f_n=0.8991469178127398.
\]

Also
\[
C(0.5)=a_{\mathrm{avg}}=454.41269746338673\ \mathrm{MeV\ fm}^3.
\]

Therefore the fixed-$y$ formulas reduce exactly to the symmetric formulas, as
required by the construction of the asymmetric fit.

\section{Summary}

The fixed-$y$ quantum SCF solver does the following:
\begin{enumerate}
\item read a fitted asymmetric Clausius interaction;
\item split the total density into proton and neutron sectors at fixed $y$;
\item compute species free-volume factors $f_p$ and $f_n$;
\item convert the physical species densities into target ideal densities;
\item solve for $\mu_p^\ast$ and $\mu_n^\ast$ through an inner damped
Newton-like SCF loop;
\item build the full pressure branch
\[
P(T,n;y)=p_p^{\mathrm{id}}+p_n^{\mathrm{id}}-\frac{C(y)n^2}{(1+cn)^2};
\]
\item solve the critical-point equations
\[
\left(\frac{\partial P}{\partial n}\right)_T=0,
\qquad
\left(\frac{\partial^2 P}{\partial n^2}\right)_T=0
\]
through an outer damped SCF loop.
\end{enumerate}

So the fixed-$y$ solver is mathematically a nested self-consistent problem:
\[
(T,n)
\Longrightarrow
(n_p,n_n)
\Longrightarrow
(f_p,f_n)
\Longrightarrow
(\mu_p^\ast,\mu_n^\ast)
\Longrightarrow
P(T,n;y)
\Longrightarrow
(T_c,n_c).
\]

\end{document}
